\documentclass[11pt,french]{article}

\usepackage[utf8]{inputenc}
\usepackage[T1]{fontenc}
\usepackage[french]{babel}
\usepackage{lmodern}
\usepackage{amsmath,amsfonts,amssymb}
\usepackage[top=2.5cm, bottom=2.5cm, left=2.5cm, right=2.5cm]{geometry}
\usepackage{booktabs}
\usepackage{array}
\usepackage{enumitem}
\usepackage{fancyhdr}
\usepackage[colorlinks=true, linkcolor=blue!70!black, urlcolor=blue!70!black]{hyperref}

\renewcommand{\baselinestretch}{1.15}
\setlength{\parindent}{0pt}
\setlength{\parskip}{0.55em}

\pagestyle{fancy}
\fancyhf{}
\fancyhead[L]{\footnotesize Université Laval \\ Département de mathématiques et de statistique}
\fancyhead[R]{\footnotesize STT-2902 -- Modélisation statistique \\ Julien Miron}
\fancyfoot[C]{\thepage}
\renewcommand{\headrulewidth}{0.5pt}
\setlength{\headheight}{28pt}

% Commandes pour les encadrés informels
\newcommand{\definition}[2]{%
  \medskip\noindent\fbox{\textbf{Définition}} \textbf{(#1)} \quad #2\medskip}
\newcommand{\remarque}[1]{%
  \medskip\noindent\textit{\textbf{Remarque.}} #1\medskip}

\begin{document}

\begin{center}
  {\LARGE\bfseries Échantillonnage}
\end{center}

\medskip
\begin{center}
  \textit{Le matériel de ce chapitre est tiré des notes de cours d'Anne-Sophie Charest et d'Emmanuelle Reny-Nolin.}
\end{center}

\bigskip
\tableofcontents

\newpage

%%============================================================
\section{Introduction}
%%============================================================

Avant de pouvoir analyser des données, il faut les collecter. La qualité de toute analyse
statistique dépend en grande partie de la façon dont les données ont été recueillies. Un
échantillon mal conçu ou entaché d'un biais peut mener à des conclusions erronées, peu
importe la rigueur des méthodes statistiques appliquées par la suite.

Ce chapitre présente les principaux concepts liés à la collecte de données~: les types
d'études statistiques, les méthodes d'échantillonnage et les sources de biais à éviter.

%%------------------------------------------------------------
\subsection{Population et échantillon}
%%------------------------------------------------------------

\textbf{Définition.} La \textbf{population} est l'ensemble de tous les individus sur lesquels
porte l'étude. Un \textbf{échantillon} est un sous-ensemble de la population qui est
effectivement observé.

En pratique, il est rarement possible d'observer toute la population (recensement). On se
contente donc d'un échantillon, à partir duquel on cherche à tirer des conclusions sur la
population entière.

\textbf{Définition.} La \textbf{base de sondage} est la liste (ou la description) de tous les
individus de la population à partir de laquelle on sélectionne l'échantillon. Idéalement,
la base de sondage correspond exactement à la population cible.

\textbf{Exemple.} Si l'on souhaite étudier la consommation de boissons énergisantes chez les
étudiants de l'Université Laval, la population cible est l'ensemble des étudiants inscrits.
La base de sondage pourrait être la liste des étudiants fournie par le bureau du registraire.

%%============================================================
\section{Types d'études}
%%============================================================

On distingue deux grands types d'études statistiques~: les \textbf{études observationnelles}
et les \textbf{expériences}.

%%------------------------------------------------------------
\subsection{Étude observationnelle}
%%------------------------------------------------------------

\textbf{Définition.} Dans une \textbf{étude observationnelle}, le chercheur observe les
individus sans intervenir sur les variables mesurées. Les individus ne sont pas assignés à
des groupes par le chercheur~: on mesure les variables telles qu'elles se présentent
naturellement.

Par conséquent, une étude observationnelle ne permet pas de conclure à un lien de
\textbf{causalité} entre deux variables. D'autres variables non mesurées, appelées
\textbf{variables confondantes}, pourraient en effet expliquer la relation observée.

\textbf{Exemple.} On observe que les étudiants de la faculté de médecine consomment davantage
de boissons énergisantes que ceux d'autres facultés. On ne peut pas conclure que le fait
d'être inscrit en médecine \emph{cause} une plus grande consommation~: il se peut que les
étudiants en médecine aient un horaire plus chargé, ce qui expliquerait les deux phénomènes
simultanément.

%%------------------------------------------------------------
\subsection{Expérience (étude expérimentale)}
%%------------------------------------------------------------

\textbf{Définition.} Dans une \textbf{expérience}, le chercheur impose délibérément un
traitement à certains individus afin d'en observer l'effet. L'\textbf{attribution aléatoire}
des individus aux groupes de traitement est la caractéristique essentielle d'une expérience
bien conçue.

Grâce à l'attribution aléatoire, les variables confondantes tendent à se répartir
uniformément entre les groupes. On peut alors attribuer les différences observées au
traitement lui-même, et ainsi conclure à un lien de causalité.

\textbf{Éléments d'une expérience bien conçue~:}
\begin{itemize}[itemsep=2pt]
  \item \textbf{Groupe traitement}~: reçoit le traitement étudié.
  \item \textbf{Groupe témoin (contrôle)}~: ne reçoit pas le traitement, ou reçoit un placebo.
  \item \textbf{Attribution aléatoire}~: les individus sont assignés aux groupes de façon
        aléatoire, indépendamment de leurs caractéristiques.
  \item \textbf{Effet placebo}~: amélioration observée chez un individu qui croit recevoir un
        traitement actif alors qu'il reçoit un placebo. C'est pourquoi les essais médicaux sont
        souvent conçus en double aveugle.
  \item \textbf{Simple aveugle}~: les sujets ne savent pas à quel groupe ils appartiennent.
  \item \textbf{Double aveugle}~: ni les sujets ni les évaluateurs ne connaissent l'appartenance
        de chaque sujet à un groupe.
\end{itemize}

\remarque{Pour pouvoir généraliser les conclusions d'une expérience à l'ensemble de la
population, les individus doivent avoir été sélectionnés de façon représentative (par exemple,
par un échantillonnage probabiliste). L'attribution aléatoire permet de conclure à la
causalité~; l'échantillonnage aléatoire permet de généraliser.}

%%============================================================
\section{Méthodes d'échantillonnage}
%%============================================================

Il existe deux grandes familles de méthodes d'échantillonnage~: les méthodes
\textbf{non probabilistes} et les méthodes \textbf{probabilistes}.

%%------------------------------------------------------------
\subsection{Échantillonnage non probabiliste}
%%------------------------------------------------------------

Dans un échantillonnage non probabiliste, les individus de la population n'ont pas tous une
probabilité connue et non nulle d'être sélectionnés. Ces méthodes sont généralement plus
faciles et moins coûteuses à mettre en oeuvre, mais elles introduisent un risque de biais et
ne permettent pas d'appliquer rigoureusement les méthodes d'inférence statistique.

\subsubsection{Échantillon de convenance}

On sélectionne les individus les plus facilement accessibles~: par exemple, demander à ses
amis ou recruter des étudiants volontaires dans un cours. Ce type d'échantillon n'est
généralement pas représentatif de la population, car les individus facilement accessibles
peuvent différer de façon systématique du reste de la population.

\subsubsection{Échantillon raisonné (ou de jugement)}

Le chercheur sélectionne lui-même les individus qu'il juge représentatifs de la population,
en se basant sur son expertise. Bien que cette méthode puisse sembler raisonnable, elle est
subjective et peut introduire des biais difficiles à détecter.

%%------------------------------------------------------------
\subsection{Échantillonnage probabiliste}
%%------------------------------------------------------------

\textbf{Définition.} Dans un \textbf{échantillonnage probabiliste}, chaque individu de la
population a une probabilité \emph{connue} et \emph{non nulle} d'être sélectionné. Ces
probabilités ne doivent pas nécessairement être égales.

L'échantillonnage probabiliste est la condition nécessaire pour utiliser les méthodes
d'inférence statistique et pour généraliser les conclusions à la population.

\subsubsection{Échantillonnage aléatoire simple (EAS)}

On sélectionne $n$ individus parmi les $N$ individus de la population de façon entièrement
aléatoire, sans remise (chaque individu a donc la même probabilité $n/N$ d'être choisi). En
pratique, on numérote les individus de la base de sondage et on utilise un générateur de
nombres aléatoires pour effectuer la sélection.

\textbf{Avantages~:} simple à comprendre et à mettre en oeuvre~; sert de référence théorique
pour les autres méthodes.

\textbf{Inconvénients~:} nécessite une base de sondage complète et à jour~; peut être
inefficace si la population est hétérogène.

\subsubsection{Échantillonnage systématique}

On choisit un individu de départ au hasard parmi les $k$ premiers, où $k = N/n$ (arrondi au
nombre entier le plus proche), puis on sélectionne tous les $k$\ieme{} individus suivants.

\textbf{Exemple.} On souhaite interroger 200 étudiants parmi 10\,000. Le pas de sélection est
$k = 10\,000/200 = 50$. On tire un nombre aléatoire $r$ entre 1 et 50, puis on sélectionne les
individus $r,\; r+50,\; r+100,\; \ldots,\; r+9\,950$.

\textbf{Avantages~:} facile à mettre en oeuvre une fois la base de sondage établie.

\textbf{Inconvénients~:} risque de biais si la liste présente une périodicité qui coïncide
avec le pas de sélection.

\subsubsection{Échantillonnage stratifié}

On divise la population en sous-groupes homogènes appelés \textbf{strates}, puis on effectue
un EAS indépendant dans chaque strate. La taille de l'échantillon dans chaque strate peut
être proportionnelle à la taille de la strate (\textbf{stratification proportionnelle}) ou
non (\textbf{stratification disproportionnelle}).

\textbf{Exemple.} Pour étudier la consommation de boissons énergisantes par faculté, on
pourrait stratifier par faculté et sélectionner un nombre d'étudiants proportionnel à la
taille de chaque faculté.

\textbf{Avantages~:} plus précis que l'EAS lorsque les strates sont homogènes en elles-mêmes
mais hétérogènes entre elles~; garantit une représentation de chaque strate dans l'échantillon.

\textbf{Inconvénients~:} nécessite de connaître l'appartenance de chaque individu à une strate
avant d'échantillonner.

\subsubsection{Échantillonnage par grappes}

On divise la population en sous-groupes appelés \textbf{grappes}, on en sélectionne quelques-unes
aléatoirement, puis on interroge tous les individus des grappes sélectionnées (plan à une phase)
ou un sous-échantillon dans chaque grappe retenue (plan à deux phases).

\textbf{Exemple.} Une université comporte 20 facultés divisées en 10 départements chacune.
On pourrait sélectionner 3 départements au hasard et interroger tous les étudiants de ces
départements.

\textbf{Avantages~:} très pratique lorsqu'il n'existe pas de liste complète de la population,
mais qu'il est facile d'identifier des groupes naturels~; réduit les coûts de collecte.

\textbf{Inconvénients~:} moins précis que l'EAS, car les individus d'une même grappe tendent
à se ressembler~; l'information apportée par chaque individu est donc partiellement redondante.

%%------------------------------------------------------------
\subsection{Résumé comparatif}
%%------------------------------------------------------------

\begin{center}
\renewcommand{\arraystretch}{1.4}
\begin{tabular}{p{3.5cm}p{5.2cm}p{5.2cm}}
\toprule
\textbf{Méthode} & \textbf{Avantages principaux} & \textbf{Inconvénients principaux} \\
\midrule
EAS & Référence théorique, équitable & Nécessite une liste complète \\
Systématique & Simple à mettre en oeuvre & Risque de biais périodique \\
Stratifié & Très précis, représentatif par strate & Nécessite la connaissance des strates \\
Par grappes & Pratique, économique & Moins précis, dépend de l'homogénéité intra-grappe \\
\bottomrule
\end{tabular}
\end{center}

%%============================================================
\section{Sources de biais}
%%============================================================

Un \textbf{biais} est une erreur \emph{systématique} qui fait que les données recueillies ne
représentent pas fidèlement la population cible. Il est important de distinguer les biais des
erreurs d'échantillonnage~: l'erreur d'échantillonnage est une variation aléatoire inhérente
à tout processus d'échantillonnage, tandis qu'un biais est une déviation constante et
orientée.

%%------------------------------------------------------------
\subsection{Biais de couverture}
%%------------------------------------------------------------

Le biais de couverture survient lorsque la base de sondage ne correspond pas exactement à la
population cible~: certains individus sont absents de la liste et ne peuvent pas être
sélectionnés.

\textbf{Exemple.} Utiliser un annuaire téléphonique de lignes fixes comme base de sondage pour
étudier la population québécoise en général exclut toutes les personnes sans ligne fixe
(notamment les jeunes adultes). Si l'âge est lié à la variable d'intérêt, cela introduit un
biais dans les estimations.

%%------------------------------------------------------------
\subsection{Biais de sélection}
%%------------------------------------------------------------

Le biais de sélection survient lorsque le mécanisme de sélection des individus favorise
systématiquement certains sous-groupes de la population. Il est particulièrement fréquent dans
les échantillonnages non probabilistes.

\textbf{Exemple.} Un sondage en ligne ne peut rejoindre que les personnes ayant accès à
Internet et habituées à ce mode de communication, ce qui exclut une partie de la population.

%%------------------------------------------------------------
\subsection{Biais de non-réponse}
%%------------------------------------------------------------

Le biais de non-réponse survient lorsque les individus qui ne répondent pas au sondage ont des
caractéristiques systématiquement différentes de ceux qui y répondent.

\remarque{Un taux de non-réponse élevé n'entraîne pas \emph{nécessairement} un biais. Un biais
apparaît seulement si les non-répondants diffèrent \emph{systématiquement} des répondants
quant à la variable d'intérêt. Cependant, un taux de non-réponse élevé augmente le risque de
biais et rend difficile la vérification de cette hypothèse.}

%%------------------------------------------------------------
\subsection{Biais de mesure}
%%------------------------------------------------------------

Le biais de mesure survient lorsque la façon de mesurer une variable -- notamment la formulation
des questions -- influence les réponses obtenues.

\textbf{Sources courantes de biais de mesure~:}
\begin{itemize}[itemsep=4pt]
  \item \textbf{Questions orientées (ou suggestives)~:} une question formulée de façon à
        suggérer une réponse.

        Formulation biaisée~: \og Êtes-vous d'accord que les boissons énergisantes sont
        nuisibles à la santé~? \fg

        Formulation neutre~: \og Quelle est votre opinion sur les boissons énergisantes~? \fg

  \item \textbf{Désirabilité sociale~:} les répondants ont tendance à donner des réponses
        socialement acceptables plutôt que leurs vraies opinions (par exemple, sous-déclarer
        leur consommation d'alcool).

  \item \textbf{Effet de l'intervieweur~:} la présence, le ton ou l'apparence de l'intervieweur
        peut influencer les réponses.

  \item \textbf{Ordre des questions~:} l'ordre dans lequel les questions sont posées peut
        influencer les réponses aux questions subséquentes.
\end{itemize}

%%============================================================
\section{Quand peut-on généraliser et conclure à la causalité~?}
%%============================================================

%%------------------------------------------------------------
\subsection{Généraliser à la population}
%%------------------------------------------------------------

On peut généraliser les résultats d'un échantillon à l'ensemble de la population lorsque
l'échantillon a été sélectionné par \textbf{échantillonnage probabiliste}. Un échantillon non
probabiliste (de convenance ou raisonné) ne permet pas de généraliser de façon rigoureuse.

%%------------------------------------------------------------
\subsection{Conclure à la causalité}
%%------------------------------------------------------------

On peut conclure à un lien de causalité entre une variable explicative et une variable réponse
uniquement lorsque les individus ont été \textbf{assignés aléatoirement} aux groupes de
traitement. Cette condition est remplie dans une expérience bien conçue.

Dans une étude observationnelle, la présence d'une association statistique entre deux variables
ne suffit pas à conclure à la causalité, car des variables confondantes pourraient être à
l'origine de la relation observée.

%%------------------------------------------------------------
\subsection{Tableau récapitulatif}
%%------------------------------------------------------------

\begin{center}
\renewcommand{\arraystretch}{1.5}
\begin{tabular}{lcc}
\toprule
\textbf{Condition} & \textbf{Généraliser à la population} & \textbf{Conclure à la causalité} \\
\midrule
Échantillonnage probabiliste uniquement & Oui & Non \\
Attribution aléatoire uniquement        & Non$^{*}$ & Oui \\
Les deux réunies                        & Oui & Oui \\
Aucune des deux                         & Non & Non \\
\bottomrule
\end{tabular}
\end{center}

{\footnotesize $^{*}$~À moins que l'expérience ait été menée sur un échantillon probabiliste
de la population.}

\bigskip

\remarque{Ces deux dimensions -- généralisation et causalité -- sont indépendantes. Une
expérience réalisée sur des étudiants volontaires d'une seule université peut permettre de
conclure à la causalité (grâce à l'attribution aléatoire), sans que ces résultats puissent
être généralisés à l'ensemble de la population (faute d'un échantillonnage probabiliste).}

\end{document}
